% Part: first-order-logic
% Chapter: models-theories
% Section: introduction

\documentclass[../../../include/open-logic-section]{subfiles}

\begin{document}

\olfileid{fol}{mat}{int}
\olsection{పరిచయం}

\begin{explain}
స్వీకృతాధార పద్ధతి అభివృద్ధి విజ్ఞానశాస్త్ర చరిత్రలో ఒక విశిష్టమైన
సాఫల్యం; గణితశాస్త్ర చరిత్రలో దానికి ప్రత్యేక ప్రాముఖ్యం ఉంది. ఒక
విషయక్షేత్రాన్ని స్వీకృతాధారంగా అభివృద్ధి చేయడంలో అనేక ప్రశ్నలను
స్పష్టం చేయవలసి వస్తుంది: ఆ క్షేత్రం దేని గురించి? అత్యంత మౌలికమైన
భావనలు ఏవి? వాటి మధ్య సంబంధం ఏమిటి? ఆ క్షేత్రంలోని భావనలన్నిటినీ ఈ
మౌలిక భావనల ద్వారా నిర్వచించగలమా? ఈ భావనలు ఏ నియమాలను పాటిస్తాయి,
ఏ నియమాలను తప్పనిసరిగా పాటించాలి?

స్వీకృతాధార పద్ధతి, తర్కం ఒకదానికొకటి సరిగ్గా సరిపోతాయి. స్వీకృతాధార
సిద్ధాంతాలను సూత్రీకరించడానికి, సిద్ధాంతపు స్వీకృతాల నుంచి ప్రమేయాలను
ఖచ్చితంగా నిర్దేశించిన విధంగా నిరూపించడానికి, ఆ స్వీకృతాలను సంతృప్తిపరిచే
వ్యవస్థలన్నిటి ధర్మాలను క్రమబద్ధంగా అధ్యయనం చేయడానికి ఔపచారిక తర్కం
పరికరాలను అందిస్తుంది.
\end{explain}

\begin{defn}
\tetoken{వాక్యాలు}{!!{sentence}s} సమితి~$\Gamma$ను, $\Gamma \Entails !A$ అయినప్పుడల్లా
$!A \in \Gamma$ అయితే, అప్పుడు మాత్రమే \emph{సంవృతం} అంటాం.
\tetoken{వాక్యాలు}{!!{sentence}s} సమితి~$\Gamma$ యొక్క \emph{సంవృతి}
$\Setabs{!A}{\Gamma \Entails !A}$.

$\Gamma$, వాక్యాల సమితి~$\Delta$ యొక్క సంవృతి అయితే, $\Delta$
ద్వారా~$\Gamma$ \emph{స్వీకృతీకరించబడింది} అంటాం.
\end{defn}

\begin{explain}
ఒక స్వీకృతాధార సిద్ధాంతాన్ని, దాని స్వీకృతాల సమితి~$\Delta$ ద్వారా
స్వీకృతీకరించబడిన వాక్యాల సమితిగా భావించవచ్చు. మరోలా చెప్పాలంటే, మనం
అధ్యయనం చేయదలచిన స్వీకృతాధార విజ్ఞానశాస్త్రపు ఆదిమ భావనలకు తార్కికేతర
సంకేతాలు గల ఒక మొదటిస్థాయి భాషతో పాటు, ఆ విజ్ఞానశాస్త్రపు మౌలిక నియమాలను
వ్యక్తం చేసే \tetoken{వాక్యాలు}{!!{sentence}s} సమితి ఉన్నప్పుడు, ఆ స్వీకృతాల నుంచి అనుగమించే ఈ
భాషలోని \tetoken{వాక్యాలు}{!!{sentence}s} అన్నిటి ద్వారా సిద్ధాంతం ప్రాతినిధ్యం పొందుతుందని
భావించవచ్చు. ఇది ఒకే ఆదిమ భావన, సరళమైన స్వీకృతాలు మాత్రమే ఉన్న పాక్షిక
క్రమాల సిద్ధాంతం వంటి సరళ ఉదాహరణల నుంచి న్యూటనీయ యాంత్రికశాస్త్రం వంటి
సంక్లిష్ట సిద్ధాంతాల వరకు విస్తరిస్తుంది.

స్వీకృతాధార పద్ధతికి ఈ ఔపచారిక దృక్పథాన్ని అంత ముఖ్యమైనదిగా చేసే ప్రధాన
తార్కిక వాస్తవాలు కిందివి. $\Gamma$ ఒక సిద్ధాంతానికి స్వీకృత వ్యవస్థ,
అంటే వాక్యాల సమితి అనుకుందాం.
\begin{enumerate}
\item ఒక స్వీకృత వ్యవస్థ ఉద్దేశించిన \tetoken{నిర్మాణాలు}{!!{structure}s} వర్గాన్ని ఎప్పుడు
  పట్టుకుంటుందో ఖచ్చితంగా చెప్పగలం. అంటే, ఒక నిర్దిష్ట \tetoken{నిర్మాణాలు}{!!{structure}s}
  వర్గంపై మనకు ఆసక్తి ఉంటే, \tetoken{నిర్మాణాలు}{!!{structure}s} సరిగ్గా
  $\Sat{M}{\Gamma}$ అయ్యే~$\Struct M$లే అయినప్పుడు మాత్రమే స్వీకృత
  వ్యవస్థ~$\Gamma$ ఆ వర్గాన్ని విజయవంతంగా పట్టుకుంటుంది.
\item $\Sat{M}{\Gamma}$ అయ్యే~$\Struct M$లు ఉన్నప్పటికీ, $\Struct M$
  మనం ఉద్దేశించిన \tetoken{నిర్మాణాలలో}{!!{structure}s} ఒకటి కాకపోవడం వల్ల ఈ విషయంలో మనం
  విఫలం కావచ్చు. అప్పుడు~$\Struct M$లో సత్యం కాని స్వీకృతాలను చేర్చాల్సి
  రావచ్చు.
\item కనీసం ఉద్దేశించిన \tetoken{నిర్మాణాలు}{!!{structure}s} అన్నిటిలో~$\Gamma$ సత్యం అనే
  విషయంలో మనం విజయవంతమైతే, $\Gamma \Entails !A$ అయినప్పుడల్లా
  వాక్యం~$!A$ ఉద్దేశించిన \tetoken{నిర్మాణాలు}{!!{structure}s} అన్నిటిలో సత్యం. కాబట్టి
  వాక్యాలు స్వీకృతాల నుంచి అనుగమిస్తాయని చూపడం ద్వారానే అవి ఉద్దేశించిన
  \tetoken{నిర్మాణాలు}{!!{structure}s} అన్నిటిలో సత్యమని చూపడానికి తార్కిక పరికరాలను
  (ఉదాహరణకు \tetoken{వ్యుత్పత్తి}{!!{derivation}} పద్ధతులను) ఉపయోగించవచ్చు.
\item కొన్నిసార్లు మన మనసులో ఉద్దేశించిన \tetoken{నిర్మాణాలు}{!!{structure}s} ఉండవు; బదులుగా
  స్వీకృతాల నుంచే మొదలుపెడతాం: కొన్ని నియమాలను సంతృప్తిపరచాలని మనం
  కోరుకునే ఆదిమ భావనలతో మొదలుపెట్టి, ఆ నియమాలను ఒక స్వీకృత వ్యవస్థలో
  సంకేతీకరిస్తాం. వెంటనే నిర్ధారించదలచే విషయం ఏమిటంటే, స్వీకృతాలు
  పరస్పరం విరుద్ధంగా లేవని. అవి విరుద్ధంగా ఉంటే ఆ నియమాలను పాటించే
  భావనలు ఏవీ ఉండవు; మనం అసంగతమైన సిద్ధాంతాన్ని ఏర్పరచడానికి
  ప్రయత్నించినట్టవుతుంది. $\Gamma$కు ఒక నమూనాను కనుగొనడం ద్వారా ఇలా
  జరగదని నిర్ధారించవచ్చు. మన సిద్ధాంతానికి నమూనాలు ఉంటే, వాటిని
  పరిశోధించడానికీ నమూనాలను నిర్మించడానికీ తార్కిక పద్ధతులను ఉపయోగించవచ్చు.
\item స్వీకృతాల స్వతంత్రత కూడా ఒక ముఖ్యమైన ప్రశ్న. స్వీకృతాల్లో ఒకటి
  నిజానికి మిగిలిన వాటి పరిణామమే అయ్యి, అందువల్ల అనవసరం కావచ్చు.
  $\Gamma$లోని స్వీకృతం~$!A$ అనవసరమని
  $\Gamma \setminus \{!A\} \Entails !A$ను నిరూపించడం ద్వారా
  చూపవచ్చు. $(\Gamma \setminus \{!A\}) \cup \{\lnot !A\}$
  సంతృప్తిపరచదగినదని చూపడం ద్వారా ఒక స్వీకృతం అనవసరం కాదని కూడా
  నిరూపించవచ్చు. ఉదాహరణకు, సమాంతర స్వీకృతం జ్యామితిలోని ఇతర
  స్వీకృతాల నుంచి స్వతంత్రమని చూపబడింది ఈ విధంగానే.
\item సిద్ధాంతంలోని భావనల నిర్వచనీయత మరో ముఖ్యమైన ప్రశ్న: భాష ఎంపిక
  సిద్ధాంతపు నమూనాలు దేనితో ఏర్పడతాయో నిర్ణయిస్తుంది. కానీ సిద్ధాంతపు
  ప్రతి అంశానికీ దాని నమూనాల్లో ప్రత్యేక ప్రాతినిధ్యం అవసరం లేదు.
  ఉదాహరణకు, ప్రతి క్రమం~$\le$ దానికి సరిపడే కఠిన క్రమం~$<$ను
  నిర్ణయిస్తుంది---వాటిలో ఒకటి ఇచ్చినప్పుడు మరొకదాన్ని నిర్వచించవచ్చు.
  కాబట్టి అలాంటి క్రమాన్ని కలిగిన సిద్ధాంతపు నమూనాలో దానికి సరిపడే కఠిన
  క్రమం \emph{కూడా} తప్పనిసరిగా ఉండనవసరం లేదు. సాధారణంగా ఒక సంబంధాన్ని
  ఇతర సంబంధాల ద్వారా ఎప్పుడు నిర్వచించవచ్చు? ఒక సంబంధాన్ని ఇతర
  సంబంధాల ద్వారా నిర్వచించడం ఎప్పుడు అసాధ్యం (అందువల్ల దానిని భాషలోని
  ఆదిమ భావనలకు చేర్చాలి)?
\end{enumerate}
\end{explain}


\end{document}
